\documentclass[aspectratio=169,10pt,english]{beamer}

\input{../../../_preambulo_beamer.tex}
\input{../../../_comandos_beamer.tex}

\selectlanguage{english}

\title{MA1001B - Statistical Modeling for Decision Making}
\subtitle{Lesson 33: Hypothesis Testing Foundations}
\author{}
\institute{Tecnol\'ogico de Monterrey}
\date{}

\begin{document}

\begin{frame}[plain]
  \vspace{1.2cm}
  {\Large\bfseries MA1001B - Statistical Modeling for Decision Making\par}
  \vspace{0.7cm}
  {\large Lesson 33: Hypothesis Testing: $H_{0}$, $H_{1}$, Type I and Type II Errors\par}
  \vspace{0.7cm}
  {\color{TecRojo}\rule{\textwidth}{1.2pt}}
  \vspace{0.8cm}

  MA1001B Faculty\\[0.3cm]
  Term\\[0.3cm]
  Tecnol\'ogico de Monterrey
\end{frame}

\begin{frame}{The Testing-Language Question}
  \begin{alertblock}{Guiding question}
    How are $H_{0}$, $H_{1}$, and $\alpha$ turned into a rejection rule, and
    what do Type I and Type II errors mean?
  \end{alertblock}

  \vspace{0.6em}
  By the end of this lesson, you can:
  \begin{itemize}
    \item state $H_{0}:\mu=50$ versus $H_{1}:\mu>50$;
    \item convert $\alpha=0.05$ into an $\bar{x}$ cutoff;
    \item interpret a Type I error as rejecting a true $H_{0}$;
    \item explain why failing to reject $H_{0}$ is not proof that $H_{0}$ is true.
  \end{itemize}

  \vfill
  \scriptsize All cycle times used today are synthetic. No real company data are used.
\end{frame}

\begin{frame}{A One-Sided Mean Test}
  \small
  \begin{center}
    \begin{tabular}{lr}
      \toprule
      \textbf{Quantity} & \textbf{Value} \\
      \midrule
      $H_{0}$ & $\mu=50$ minutes \\
      $H_{1}$ & $\mu>50$ minutes \\
      $n$, known $\sigma$ & 36, 8 minutes \\
      $\alpha$ & 0.05 \\
      \bottomrule
    \end{tabular}
  \end{center}

  \vspace{0.5em}
  \begin{exampleblock}{Cutoff}
    \[
      \mathrm{SE}=\frac{8}{\sqrt{36}}=1.333333,\qquad
      z^{*}=1.644854,
    \]
    \[
      \bar{x}_{\text{crit}}=50+1.644854\times 1.333333=52.193138.
    \]
  \end{exampleblock}
\end{frame}

\begin{frame}[fragile]{Step 1: Hypotheses and Alpha}
  \begin{columns}[T]
    \begin{column}{0.42\textwidth}
      \small
      \begin{lstlisting}
z_crit = stats.norm.ppf(0.95)
xbar_crit = (
    mu0 + z_crit * se
)
      \end{lstlisting}

      \begin{block}{Verified output}
        $z^{*}=1.644854$\\
        Cutoff $=52.193138$
      \end{block}
    \end{column}
    \begin{column}{0.58\textwidth}
      \centering
      \includegraphics[width=\textwidth]{../figures/ht_foundations_01_hypotheses_alpha.png}
      \vspace{0.2em}
      \scriptsize Rejection region from Step 1
    \end{column}
  \end{columns}
\end{frame}

\begin{frame}{Type I Error: Rejecting a True $H_{0}$}
  \small
  A Type I error occurs when $H_{0}$ is true and the sample still lands in
  the rejection region. The long-run rate of that event is $\alpha$.

  \vspace{0.5em}
  Seed-42 simulation of $10{,}000$ samples from $\mu=50$:
  \[
    531 \text{ rejections},\qquad
    \widehat{\alpha}=0.053100.
  \]
  The simulated rate sits near the nominal $0.05$. It is not identical to
  $0.05$ in a finite Monte Carlo.
\end{frame}

\begin{frame}[fragile]{Step 2: Simulated Type I Rate}
  \begin{columns}[T]
    \begin{column}{0.40\textwidth}
      \small
      \begin{lstlisting}
xbars = samples.mean(axis=1)
rate = np.mean(
    xbars >= xbar_crit
)
      \end{lstlisting}

      \begin{block}{Verified output}
        Rejections $=531$\\
        Type I rate $=0.053100$
      \end{block}
    \end{column}
    \begin{column}{0.60\textwidth}
      \centering
      \includegraphics[width=\textwidth]{../figures/ht_foundations_02_type_i_rate.png}
      \vspace{0.2em}
      \scriptsize $H_{0}$ sampling distribution from Step 2
    \end{column}
  \end{columns}
\end{frame}

\begin{frame}{Step 3: Failing to Reject Is Not Proof}
  \begin{columns}[T]
    \begin{column}{0.42\textwidth}
      \small
      Power at $\mu=53$: $0.724100$.

      \vspace{0.4em}
      Type II rate: $0.275900$.

      \vspace{0.4em}
      Nearby $\mu=51$ sample:
      $\bar{x}=51.590798<52.193138$.
      $H_{0}$ is not rejected.

      \vspace{0.5em}
      \textbf{Limit:} $H_{1}$ can be true and the test can still fail to
      reject. That non-rejection is not proof that $\mu=50$.
    \end{column}
    \begin{column}{0.58\textwidth}
      \centering
      \includegraphics[width=\textwidth]{../figures/ht_foundations_03_type_ii_not_proof.png}
      \vspace{0.2em}
      \scriptsize Power and Type II from Step 3
    \end{column}
  \end{columns}
\end{frame}

\begin{frame}{Concept Check}
  \begin{enumerate}
    \item Write $H_{0}$ and $H_{1}$ for this packing-time test.
    \item What is a Type I error in words?
    \item Why is a simulated Type I rate of $0.053100$ acceptable here?
    \item A nearby $\mu=51$ sample gives $\bar{x}=51.59$. Why is that not
      proof that $\mu=50$?
  \end{enumerate}

  \vspace{0.6em}
  \begin{block}{Connection to Lesson 34}
    The session can continue directly into the one-sample z test: the $z$
    statistic, the $p$-value, and the known-sigma assumption.
  \end{block}

  \vfill
  \scriptsize Source: OpenStax, \textit{Introductory Business Statistics 2e},
  Chapter 9, Sections 9.1--9.2. CC BY-NC-SA.
\end{frame}

\appendix



\end{document}
